\documentclass[11pt,a4paper]{article}

% Traxia-style arXiv preamble (pdfLaTeX; no fontspec)
\usepackage[english]{babel}
\usepackage[utf8]{inputenc}
\usepackage[T1]{fontenc}
\usepackage{lmodern}
\usepackage{microtype}

\usepackage{amsmath,amssymb}
\usepackage{booktabs}
\usepackage{array}
\usepackage{ragged2e}
\usepackage{graphicx}
\usepackage{caption}
\usepackage{enumitem}
\usepackage{setspace}
\usepackage[dvipsnames]{xcolor}

\usepackage[a4paper]{geometry}
\geometry{
  left=2.4cm,
  right=2.4cm,
  top=2.5cm,
  bottom=2.6cm,
  footskip=28pt,
}

\setstretch{1.05}
\setlength{\emergencystretch}{2.5em}
\setlength{\parskip}{0.3em plus 0.08em minus 0.15em}
\setlength{\parindent}{1.25em}
\setlist[itemize]{noitemsep, topsep=2pt}
\setlist[enumerate]{noitemsep, topsep=2pt}

\captionsetup{
  font=small,
  labelfont=bf,
  labelsep=colon,
  skip=6pt,
  justification=justified,
  singlelinecheck=false,
}
\captionsetup[table]{position=above, skip=5pt}

\setcounter{secnumdepth}{3}
\usepackage{titlesec}
\titleformat{\section}{\large\bfseries}{\thesection}{1em}{}
\titleformat{\subsection}{\normalsize\bfseries}{\thesubsection}{1em}{}

\renewcommand{\abstractname}{Abstract}
\renewenvironment{abstract}{%
  \centerline{\textbf{\abstractname}}\par\vspace{0.5em}%
  \small\justifying
}{%
  \par\vspace{0.6em}%
}

\usepackage[numbers,sort,compress]{natbib}
\usepackage{url}
\urlstyle{same}
\usepackage[
  colorlinks=true,
  linkcolor=NavyBlue,
  citecolor=NavyBlue,
  urlcolor=NavyBlue,
  filecolor=NavyBlue,
  bookmarks=true,
  bookmarksopen=true,
  pdfborder={0 0 0},
  breaklinks=true,
]{hyperref}
\usepackage{bookmark}
\hypersetup{
  pdftitle={Tool-Augmented Language Model Agents for Explainable Mine Safety Risk Analysis},
  pdfauthor={Wisdom Dogah},
}

\newcommand{\email}[1]{\href{mailto:#1}{#1}}

\title{%
  \textbf{Tool-Augmented Language Model Agents for}\\[0.35em]
  \textbf{Explainable Mine Safety Risk Analysis:}\\[0.35em]
  \textbf{A Study Using U.S.\ Mine Safety and Health Administration Data}%
}

\author{%
  Wisdom Dogah%
  \thanks{\footnotesize Faculty of Computing and Mathematical Sciences,\\ University of Mines and Technology (UMaT), Tarkwa, Ghana.}%
  \thanks{\footnotesize BlackMatrix AI Research, Accra, Ghana.}%
  \thanks{\footnotesize Correspondence: \email{wisdom@blackmatrix.io}.}%
  \thanks{\footnotesize Preprint. Comments welcome.}%
}

\date{July 2026}

\begin{document}
\maketitle
\thispagestyle{plain}
\justifying

\begin{abstract}
Mining remains one of the more hazardous industrial occupations, and
safety regulators such as the United States Mine Safety and Health
Administration (MSHA) collect large volumes of accident and injury
data every year, combining structured fields with free text narratives
written by mine operators. Prior work on this dataset has treated it as
a supervised learning problem: predicting injury severity or days away
from work using logistic regression, decision trees, or neural networks
trained on structured fields and, in some cases, the narrative text.
These models produce a prediction but not an explanation a safety
officer can act on, and none of them let a user ask a follow-up question
in plain language. Separately, a recent line of work on tool-augmented
large language model agents has shown that an LLM can plan, call
external tools, and reason over heterogeneous operational data in
domains such as oil and gas drilling, but this pattern has not yet been
applied to occupational mine safety data, and none of the agentic systems
in adjacent domains report a human evaluation of the explanations they
produce. This paper proposes and evaluates a tool-augmented LLM agent
that answers natural language safety questions by combining a structured
risk classifier, a trend analysis tool, and a retrieval tool over
historical injury narratives from the MSHA Accident, Injury, and Illness
dataset. The system is compared against a plain classifier baseline and
a retrieval-augmented generation baseline on answer accuracy, tool
selection correctness, and latency. On a 60-question benchmark (20
classification, 20 trend, 20 case-grounded), the primary evaluation uses a
live LLM (Llama 3.3 70B via Groq) that selects tools from natural language.
The tool-augmented agent achieves \textbf{38.3\% overall accuracy} versus
\textbf{30.0\%} for the classifier baseline and \textbf{28.3\%} for the RAG
baseline, with 60\% tool selection correctness. All systems attempt every
question. A deterministic offline router (no LLM reasoning) reaches
\textbf{93.3\%} accuracy as a zero-cost ablation ceiling only.
The structured classifier tool alone achieves 0.574 accuracy and 0.562
macro F1 on a 48{,}128-record holdout over ten injury severity classes.
Explanation quality assessment using the Explanation Satisfaction Scale
of Hoffman et al.~\cite{Hoffman2023} is prepared for mining
engineering faculty and senior students; participant data collection
is ongoing.

\noindent\textbf{Keywords:}
mine safety;
MSHA;
tool-augmented LLM;
agentic AI;
explainable AI;
occupational safety;
retrieval-augmented generation;
human evaluation.
\end{abstract}

\section{Introduction}

Mining safety has improved substantially over the past several decades
through regulation, mechanization, and safety culture, but injuries and
fatalities have not been eliminated, and mine operators and safety
officers still spend considerable effort reviewing historical incident
records by hand or through simple dashboards when trying to understand
what drives risk at a given site. The United States Mine Safety and
Health Administration has published detailed accident, injury, and
illness records since 2000 under its Open Government Data Initiative, and
this dataset has become a standard resource for applied machine learning
research on occupational mine safety. Existing studies using this dataset
fall into a narrow methodological pattern: a model is trained to predict
an outcome, such as degree of injury or number of days away from work,
from structured fields and sometimes from the injury narrative text, and
its predictive accuracy is reported.

This pattern has two limitations that matter for how the data is actually
used in practice. First, a single number or classification label does not
tell a safety officer why the model produced that answer, and does not
let them ask a natural follow-up question such as which occupations show
the sharpest recent increase in a given injury type, or whether a specific
incident resembles a known historical cluster. Second, none of the
published work on this dataset evaluates whether the output is actually
useful or trustworthy to the people who would use it, since none of it
involves a human evaluation study.

At the same time, a separate and recent strand of research has demonstrated
that large language models can be wrapped in an agent loop that plans
which tool to call, executes that tool, and reasons over the result, rather
than simply generating a single answer from a prompt. This pattern, sometimes
called tool use or tool augmentation, was formalized in general purpose form
by Yao et al.~\cite{Yao2023} and Schick et al.~\cite{Schick2023}, and has
recently been applied to a real industrial operations problem by Lu~\cite{Lu2026},
whose TADI system orchestrates twelve domain-specific tools over drilling
reports and sensor data from an oil field, explicitly avoiding heavyweight
agent frameworks in favor of a small, auditable orchestration loop built
directly on a large language model provider's function calling interface.
As far as we can determine from a review of the mining, safety, and agentic
AI literature, this pattern has not yet been applied to mine safety data,
and the general MLLM agent literature that does touch on mining, notably
the MineAgent system of Wang et al.~\cite{Wang2024}, addresses mineral
exploration from remote sensing imagery rather than occupational safety or
operational decision support.

This paper makes three contributions. First, it introduces a tool-augmented
LLM agent architecture for reasoning over the MSHA accident and injury
dataset, combining a structured classifier, a statistical trend tool, and
a retrieval tool over incident narratives, following the transparent,
framework-free orchestration design demonstrated by Lu~\cite{Lu2026} in a
different industrial domain. Second, it evaluates this system against two
baselines, a direct classifier and a single-shot retrieval-augmented
generation pipeline, on task accuracy, tool selection correctness, and
latency on a fixed 60-question benchmark. Third, it provides a human
evaluation protocol based on the Explanation Satisfaction Scale of Hoffman
et al.~\cite{Hoffman2023} for mining engineering faculty and senior
students; participant data collection is in progress.

The remainder of this paper is organized as follows. Section~\ref{sec:related}
reviews related work. Section~\ref{sec:data} describes the dataset and
preprocessing. Section~\ref{sec:architecture} describes the proposed
architecture. Section~\ref{sec:experiments} describes the experimental
design. Section~\ref{sec:results} reports results. Section~\ref{sec:discussion}
discusses contributions and limitations. Implementation steps are documented
in the repository reproduction guide.

\section{Related Work}
\label{sec:related}

\subsection{Predictive modeling of MSHA accident and injury data}

The MSHA Accident, Injury, and Illness dataset, collected under 30 CFR Part
50 and published through MSHA's Open Government Data Initiative~\cite{MSHA2024},
has been used in several prior studies. Amoako et al.~\cite{Amoako2021} applied
multiclass logistic regression to a ten-year injury dataset to identify risk
factors associated with different injury classes. Yedla et al.~\cite{Yedla2020}
used roughly 228{,}000 records with fifty variables to train deep neural
networks predicting both injury outcome and days away from work, comparing
against logistic regression and exploring synthetic data augmentation on
narrative text. Yedla~\cite{Yedla2019} applied decision trees, random forests,
and artificial neural networks to accidents reported between 2000 and 2018.
Across all three studies, the narrative text field, when used at all, is treated
as an additional input feature for a supervised model rather than as something
a user can query directly, and none includes a human evaluation of how useful
or interpretable practitioners find the output.

\subsection{Computer vision, digital twins, and multimodal AI in mining safety}

A parallel body of applied work uses computer vision to detect hazards,
personal protective equipment violations, and unsafe proximity between workers
and equipment from camera feeds. Several of these papers explicitly identify
explainable AI and multimodal sensor fusion as future work rather than
delivering it, which is consistent with the gap this paper addresses, though
these systems operate on a different data modality (real-time video) than the
one used here. Digital twin approaches to mining operations have been the
subject of recent systematic reviews; these indicate that the digital twin
literature is comparatively mature and would need a narrow, specific technical
angle to contribute something new, which is part of why this paper does not
pursue that direction.

\subsection{Tool-augmented and agentic LLM systems}

The general capability of large language models to interleave reasoning with
calls to external tools was formalized by Yao et al.~\cite{Yao2023} in the
ReAct framework and extended by Schick et al.~\cite{Schick2023} in Toolformer.
Lu~\cite{Lu2026} applied this pattern to structured drilling records and free
text daily drilling reports from the Equinor Volve field, orchestrating twelve
domain-specific tools through direct function calling. Wang et al.~\cite{Wang2024}
apply a multimodal LLM agent to mineral exploration from remote sensing
imagery, which shares a domain label with the present work but addresses a
different task rather than operational or safety decision support. No published
work identified in this review applies the tool-augmented LLM agent pattern
to occupational safety data in mining.

\subsection{Explainability and trust evaluation}

Hoffman et al.~\cite{Hoffman2023} developed validated scales for explainable
AI research covering explanation goodness, user satisfaction, curiosity, and
trust, including the Explanation Satisfaction Scale. Trust in automation more
broadly has a long history in human factors research~\cite{Lee2004}, whose
framework for appropriate reliance underlies much of the more recent human-AI
collaboration literature this project draws on for study design.

\section{Data}
\label{sec:data}

The primary dataset is the MSHA Accident, Injury, and Illness dataset,
available through MSHA's Open Government Data Initiative~\cite{MSHA2024},
containing all accidents, injuries, and illnesses reported by mine operators
and contractors in the United States since January 1, 2000, drawn from MSHA
Form 7000-1. Each record includes a unique document number and structured
fields describing subunit, degree of injury, mining equipment involved,
accident classification, occupation, activity at the time of the accident,
injury source, nature of injury, and body part affected, along with a free
text narrative describing the incident in the operator's own words.

We downloaded the Accident Injuries dataset and the Mine Identification
dataset from MSHA's data portal, removed records with missing or invalid
values in fields required for classification, and held out a stratified 20\%
test sample across calendar years and injury classes. The final cleaned corpus
contains 240{,}640 records (192{,}512 train / 48{,}128 test). Narrative text
was embedded with \texttt{sentence-transformers/all-MiniLM-L6-v2} and indexed
in ChromaDB for semantic retrieval. A benchmark of 60 natural language
questions, written before system evaluation, spans classification, trend, and
case-grounded categories (Section~\ref{sec:experiments}). Because the dataset
covers United States operations only, findings may not transfer directly to
other regulatory environments, including Ghana, without further validation.

\section{System Architecture}
\label{sec:architecture}

The system follows the general orchestration pattern demonstrated by Lu~\cite{Lu2026}
for drilling operations, adapted here for occupational safety reasoning rather
than drilling engineering, and built without a heavyweight agent framework so
that every step in the reasoning process is visible and can be audited.

\textbf{Data layer.} Structured MSHA fields are loaded into a relational
analytical store. Narrative texts are chunked and embedded into a vector store
to support semantic retrieval over historical incident descriptions.

\textbf{Tool layer.} Three tools are implemented as deterministic, testable
functions:
\begin{itemize}
  \item A \emph{risk classification tool}, which reproduces and extends the
    supervised modeling approach of prior MSHA studies (a random forest
    baseline) to predict injury severity class for a specified set of
    conditions.
  \item A \emph{trend analysis tool}, which computes rates, counts, and changes
    over time for a specified injury type, occupation, equipment category, or
    mine site.
  \item A \emph{narrative retrieval tool}, which performs semantic search over
    historical incident narratives and returns the most relevant matching
    incidents with their structured metadata.
\end{itemize}

\textbf{Orchestrator.} A large language model receives the user's question,
a system prompt describing the three tools, and access to the tools through
the provider's native function calling interface. The model selects which tool
or sequence of tools to call, receives the tool output, and produces a natural
language answer citing which tools were used.

\textbf{Baselines for comparison.} Two baselines isolate the value of tool
augmentation: a direct classifier baseline, which answers only classification-style
questions; and a single-shot retrieval-augmented generation baseline, which
retrieves relevant narrative chunks and passes them to the language model in
a single prompt without tool orchestration or access to the trend or
classification tools.

\section{Experimental Design}
\label{sec:experiments}

\textbf{Research questions.} Does tool-augmented orchestration improve answer
accuracy over a direct classifier and over a single-shot retrieval-augmented
generation baseline, for natural language safety questions posed against the
MSHA dataset? Does the system correctly select the tool or tools appropriate
to a given question? What is the latency per query, and how does this scale
with question complexity? Do mining engineering domain experts rate the
system's explanations as satisfactory and trustworthy using a validated
instrument?

\textbf{Benchmark construction.} Sixty natural language questions were authored
in three equal categories: classification-style questions with explicit MSHA
field codes, trend and comparison questions requiring aggregation, and
case-grounded questions requiring retrieval of similar historical incidents.
Reference answers were derived from cleaned data and tool outputs before any
benchmark run.

\textbf{Quantitative metrics.} Task accuracy (proportion correct per category
and overall), tool selection correctness (invoked tools vs.\ expected tools
per question), and latency (wall-clock time per query).

\textbf{Human evaluation.} Mining engineering faculty and senior students at
the University of Mines and Technology will rate blinded question--answer pairs
using the Explanation Satisfaction Scale~\cite{Hoffman2023}. Target sample:
10--20 participants; results will be reported descriptively as an exploratory
study.

\textbf{Failure case analysis.} Incorrect benchmark answers were logged and
categorized by failure type (classifier error, retrieval miss, tool routing
error).

\section{Results}
\label{sec:results}

All quantitative results were produced from the open source implementation.
Benchmark questions and reference answers were fixed before evaluation. The
primary agent evaluation used deterministic tool routing without LLM inference,
isolating tool correctness at zero API cost. Supplementary live LLM runs
(local Ollama) are reported in Section~\ref{sec:live-llm}.

\subsection{Data and classifier tool}

After cleaning, \textbf{240{,}640} accident records remained (273{,}614 raw;
exclusions documented in the reproduction guide). An 80/20 stratified split
yielded 192{,}512 training and 48{,}128 test records across ten
\texttt{DEGREE\_INJURY\_CD} classes (codes 01--10).

The random forest classifier (\texttt{InjuryRiskClassifier}, 100 trees,
balanced class weights) achieved on the stratified holdout:

\begin{table}[htbp]
  \centering
  \caption{Classifier holdout metrics on 48{,}128 test records.}
  \label{tab:classifier}
  \begin{tabular}{@{}lc@{}}
    \toprule
    Metric & Value \\
    \midrule
    Accuracy & 0.574 \\
    Macro F1 & 0.562 \\
    Weighted F1 & 0.565 \\
    Fatality (01) recall & 0.538 \\
    \bottomrule
  \end{tabular}
\end{table}

Out-of-time evaluation (train 2000--2020, test 2021+) yielded accuracy 0.553
and macro F1 0.559. These figures are comparable in spirit to prior MSHA
supervised modeling at similar data scale~\cite{Yedla2020} but are not directly
comparable across different targets and feature sets.

\subsection{Primary benchmark (Groq llama-3.3-70b)}

The primary reported result uses Groq \texttt{llama-3.3-70b-versatile}. The LLM
plans tool calls from natural language. All three systems attempt all 60
questions.

\begin{table}[htbp]
  \centering
  \caption{Primary live LLM benchmark (Groq).}
  \label{tab:groq}
  \begin{tabular}{@{}lccccc@{}}
    \toprule
    System & Class. & Trend & Case & Overall & Tool sel. \\
    \midrule
    Tool-augmented agent & 85.0\% & 5.0\% & 25.0\% & \textbf{38.3\%} & 60.0\% \\
    Classifier baseline & 90.0\% & 0.0\% & 0.0\% & 30.0\% & 100\% \\
    RAG baseline & 0.0\% & 0.0\% & 85.0\% & 28.3\% & 98.3\% \\
    \bottomrule
  \end{tabular}
\end{table}

Mean latency: agent 126\,s, classifier 0.19\,s, RAG 256\,s per question.

\subsection{Offline routing ablation (deterministic, no LLM)}

A separate run uses category metadata to route questions without LLM inference.
This measures a zero-cost routing ceiling, not live agentic reasoning.

\begin{table}[htbp]
  \centering
  \caption{Offline routing ablation ($n=20$ per category).}
  \label{tab:benchmark}
  \begin{tabular}{@{}lcccc@{}}
    \toprule
    System & Classification & Trend & Case-grounded & Overall \\
    \midrule
    Offline router & 90.0\% & \textbf{100.0\%} & 90.0\% & \textbf{93.3\%} \\
    Classifier baseline & 90.0\% & 0.0\% & 0.0\% & 30.0\% \\
    RAG baseline (no LLM) & 0.0\% & 0.0\% & 90.0\% & 30.0\% \\
    \bottomrule
  \end{tabular}
\end{table}

The 30\% overall scores for single-tool baselines reflect incorrect answers on
out-of-strength questions after attempting every question, not category blocking.
Mean latency: offline router 0.31\,s, classifier 0.22\,s, RAG 0.03\,s.

\subsection{Failure analysis (offline router)}

Four agent failures occurred:
\begin{enumerate}
  \item \textbf{Two classification errors (CLS-03, CLS-14):} classifier
    predicted degree code 03 instead of the reference code.
  \item \textbf{Two retrieval errors (CASE-14, CASE-15):} semantically similar
    narratives were retrieved but the reference document was not in the top five
    results.
\end{enumerate}
Baseline failures on out-of-domain question types are expected when a single
tool is used for every question.

\subsection{Human evaluation}

Materials for the Hoffman et al.~\cite{Hoffman2023} Explanation Satisfaction
Scale are prepared. Participant ratings are not yet included; data collection at
the University of Mines and Technology may proceed now that the primary Groq
benchmark is complete.

\subsection{Live LLM evaluation (Ollama, supplementary)}
\label{sec:live-llm}

A supplementary run used Ollama with \texttt{qwen2.5:7b} before baseline
correction (agent 53.3\% overall, 88.3\% tool selection). That run is not
directly comparable to the corrected Groq comparison above.

\section{Contributions and Limitations}
\label{sec:discussion}

This paper makes three contributions, now supported by implemented code and
measured results. First, it introduces a tool-augmented LLM agent architecture
for reasoning over the MSHA accident and injury dataset. Second, on a fixed
60-question benchmark where every system attempts every question, the live Groq
agent achieves 38.3\% overall accuracy versus 30.0\% for the classifier baseline
and 28.3\% for the RAG baseline. Third, it prepares the first human evaluation
protocol for explanation quality in mine safety AI using the Explanation
Satisfaction Scale of Hoffman et al.~\cite{Hoffman2023}; participant data
collection remains future work.

Limitations stated plainly:
\begin{itemize}
  \item \textbf{Baseline comparison correction:} An earlier draft restricted
    each baseline to its home category; the corrected comparison attempts all
    60 questions with all systems.
  \item \textbf{Offline routing ablation:} The 93.3\% offline router result uses
    category metadata, not live LLM reasoning.
  \item \textbf{Live LLM weakness:} Primary Groq agent accuracy (38.3\%) is far
    below the offline ablation; trend questions scored 5\%.
  \item \textbf{Geographic scope:} U.S.\ MSHA data only; findings may not transfer
    to Ghana or other regulatory environments without validation.
  \item \textbf{Classifier weakness:} Macro F1 0.562; several severity classes
    (04, 09, 10) have low recall.
  \item \textbf{Retrieval:} Semantic search sometimes returns plausible but
    incorrect reference documents (two offline router failures).
  \item \textbf{Human evaluation:} Protocol prepared; participant data not yet
    collected.
\end{itemize}

\section*{Acknowledgments}

We thank the University of Mines and Technology for institutional support and
the MSHA Open Government Data Initiative for public access to accident and
injury records.

\bibliographystyle{plainnat}
\bibliography{references}

\end{document}
